\begin{proof}
    The proof follows by induction.
    \begin{itemize}
    \item \textbf{Base Case}.
          Assume that $g \in \Sigma_{\Core}^0$.
          There are six cases to consider.
          \begin{enumerate}
          \item Assume that $g = \scalebox{0.65}{\circimg{tikz-cache/encoder__eqn6_lhs.pdf}}$.
                Then by definition $\scalebox{0.65}{\circimg{tikz-cache/encoder__eqn6_lhs.pdf}} \in \decoder( \scalebox{0.65}{\circimg{tikz-cache/encoder__eqn6_lhs.pdf}} ) = \decoder( \encoder( \scalebox{0.65}{\circimg{tikz-cache/encoder__eqn6_lhs.pdf}} ) )$.
                Since $( \approx_E )$ is reflexive and $g \in \decoder( \encoder( g ) )$, then there exists some $U \in \decoder( \encoder( g ) )$ such that $g \approx_E U$.
          \item Assume that $g = \scalebox{0.65}{\circimg{tikz-cache/encoder__eqn3_lhs.pdf}}$.
                Then by definition $\scalebox{0.65}{\circimg{tikz-cache/proofs__round_c3_1.pdf}} \in \decoder( \scalebox{0.65}{\circimg{tikz-cache/encoder__eqn3_rhs.pdf}} ) = \decoder( \encoder( \scalebox{0.65}{\circimg{tikz-cache/encoder__eqn3_lhs.pdf}} ) )$.
                Since $\alpha\pi/\pi = \pi$, then $g  \in \decoder( \encoder( g ) )$.
                Since $( \approx_E )$ is also reflexive, then there exists some $U \in \decoder( \encoder( g ) )$ such that $g \approx_E U$.
          \item Assume that $g = \scalebox{0.65}{\circimg{tikz-cache/encoder__eqn1_lhs.pdf}}$.
                Then the following equation holds  (see~\cref{App:CNOT}).
                \begin{align*}
                    \scalebox{0.65}{\circimg{tikz-cache/proofs__round_c1_1.pdf}}
                    \in
                    \decoder\left(
                        \scalebox{0.65}{\circimg{tikz-cache/encoder__eqn1_rhs.pdf}}
                    \right)
                    =
                    \decoder\left( \encoder\left(
                        \scalebox{0.65}{\circimg{tikz-cache/encoder__eqn1_lhs.pdf}}
                    \right) \right)
                \end{align*}
                Then the following derivation holds.
                \begin{align*}
                    \scalebox{0.65}{\circimg{tikz-cache/proofs__round_c1_1.pdf}}
                    =
                    \scalebox{0.65}{\circimg{tikz-cache/proofs__round_c1_2.pdf}}
                    \overset{\HqcFreePhase}{\approx}
                    \scalebox{0.65}{\circimg{tikz-cache/proofs__round_c1_3.pdf}}
                \end{align*}
          \item Assume that $g = \scalebox{0.65}{\circimg{tikz-cache/encoder__eqn4_lhs.pdf}}$.
                It follows from case (3) that following equation holds.
                \begin{align*}
                    \scalebox{0.65}{\circimg{tikz-cache/proofs__round_c4_1.pdf}}
                    \in
                    \decoder\left(
                        \scalebox{0.65}{\circimg{tikz-cache/encoder__eqn4_rhs.pdf}}
                    \right)
                    =
                    \decoder\left( \encoder\left(
                        \scalebox{0.65}{\circimg{tikz-cache/encoder__eqn4_lhs.pdf}}
                    \right) \right)
                \end{align*}
                Then the following derivation holds.
                \begin{align*}
                    \scalebox{0.65}{\circimg{tikz-cache/proofs__round_c4_1.pdf}}
                    \overset{\HqcPowDag}{\approx}
                    \scalebox{0.65}{\circimg{tikz-cache/proofs__round_c4_2.pdf}}
                    \overunderset{\HqcPowCong}{\HqcTargDef}{\approx}
                    \scalebox{0.65}{\circimg{tikz-cache/encoder__eqn4_lhs.pdf}}
                \end{align*}
          \item Assume that $g = \scalebox{0.65}{\circimg{tikz-cache/encoder__eqn5_lhs.pdf}}$.
                It follows from cases (3--4) that following equation holds.
                \begin{align*}
                    \scalebox{0.65}{\circimg{tikz-cache/proofs__round_c5_1.pdf}}
                    \in
                    \decoder\left(
                        \scalebox{0.65}{\circimg{tikz-cache/encoder__eqn5_rhs.pdf}}
                    \right)
                    =
                    \decoder\left( \encoder\left(
                        \scalebox{0.65}{\circimg{tikz-cache/encoder__eqn5_lhs.pdf}}
                    \right) \right)
                \end{align*}
                Then the following derivation holds.
                \begin{align*}
                    \scalebox{0.65}{\circimg{tikz-cache/proofs__round_c5_1.pdf}}
                    \overset{\HqcPowId}{\approx}
                    \scalebox{0.65}{\circimg{tikz-cache/proofs__round_c5_2.pdf}}
                    \overset{\HqcPowId}{\approx}
                    \scalebox{0.65}{\circimg{tikz-cache/proofs__round_c5_3.pdf}}
                    \overset{\HqcPowDag}{\approx}
                    \scalebox{0.65}{\circimg{tikz-cache/proofs__round_c5_4.pdf}}
                    \overunderset{\HqcPowCong}{\HqcTargMDef}{\approx}
                    \scalebox{0.65}{\circimg{tikz-cache/encoder__eqn5_lhs.pdf}}
                \end{align*}
          \item Assume that $g = \scalebox{0.65}{\circimg{tikz-cache/encoder__eqn2_lhs.pdf}}$.
                It follows from cases (3--4) that following equation holds.
                \begin{align*}
                    \scalebox{0.65}{\circimg{tikz-cache/proofs__round_c2_1.pdf}}
                    \in
                    \decoder\left(
                        \scalebox{0.65}{\circimg{tikz-cache/encoder__eqn2_rhs.pdf}}
                    \right)
                    =
                    \decoder\left( \encoder\left(
                        \scalebox{0.65}{\circimg{tikz-cache/encoder__eqn2_lhs.pdf}}
                    \right) \right)
                \end{align*}
                Then the following derivation holds.
                \begin{align*}
                    \scalebox{0.65}{\circimg{tikz-cache/proofs__round_c2_1.pdf}}
                    \overset{\HqcPowId}{\approx}
                    \scalebox{0.65}{\circimg{tikz-cache/proofs__round_c2_2.pdf}}
                    \overset{\HqcPowId}{\approx}
                    \scalebox{0.65}{\circimg{tikz-cache/proofs__round_c2_3.pdf}}
                    \overset{\HqcPowDag}{\approx}
                    \scalebox{0.65}{\circimg{tikz-cache/proofs__round_c2_4.pdf}}
                    \overunderset{\HqcPowCong}{\HqcTargMDef}{\approx}
                    \scalebox{0.65}{\circimg{tikz-cache/encoder__eqn2_lhs.pdf}}
                \end{align*}
          \end{enumerate}
          In each case, there exists a $U \in \decoder( \encoder( g ) )$ such that $g \approx_E U$.
    \item \textbf{Inductive Hypothesis}.
          For each $n \in \mathbb{N}$, if $g \in \Sigma_{\Core}^n$, then there exists some circuit $U \in \decoder( \encoder( g ) )$ such that $g \approx_E U$.
    \item \textbf{Inductive Step}.
          Assume that the inductive hypothesis holds for some $n \in \mathbb{N}$ and let $g \in \Sigma_{\Core}^{n+1}$.
          If $g \in \Sigma_{\Core}^n$ as well, then there exists some $U \in \decoder( \encoder( U ) )$ such that $g \approx_E U$ by the inductive hypothesis.
          Assume instead that $g \in \Sigma_{\Core}^{n+1} \setminus \Sigma_{\Core}^n$.
          Then there exists some $V \in \Mor( \Core[n] )$ such that $g = \scalebox{0.65}{\circimg{tikz-cache/dagger__term1.pdf}} \odot V$.
          Then by the definitions of $\encoder( - )$ and $\decoder( - )$, the following equation holds.
          \begin{equation*}
              \decoder( \encoder( g ) )
              =
              \decoder( C( \encoder( V ) ) )
              =
              C( \decoder( \encoder( V ) ) )
          \end{equation*}
          Then by the inductive hypothesis, there exists some $W \in \decoder( \encoder( V ) )$ such that $V \approx_E W$.
          Then the following equations hold  (see~\cref{App:CNOT}).
          \begin{equation*}
              \scalebox{0.65}{\circimg{tikz-cache/proofs__round_ind_1.pdf}}
              \in
              \decoder\left( \encoder\left(
                \scalebox{0.65}{\circimg{tikz-cache/proofs__round_ind_2.pdf}}
              \right) \right)
              %
              \qquad\qquad
              %
              \scalebox{0.65}{\circimg{tikz-cache/proofs__round_ind_1.pdf}}
              \overunderset{\HQCCtrlCong}{I.H.}{\approx}
              \scalebox{0.65}{\circimg{tikz-cache/proofs__round_ind_2.pdf}}
          \end{equation*}
          Then there exists some $U \in \decoder( \encoder( g ) )$ such that $g \approx_E U$
    \end{itemize}
    Then by the principle of inductive, for each $g \in \Sigma_{\Core}$, there exists some $U \in \decoder( \encoder( g ) )$ such that $g \approx_E U$.
\end{proof}
